\section{The Seiberg-Witten Analysis of $N=2$ Supersymmetric
Yang-Mills Theory}

In this section, we analyze the $N = 2$ pure $SU(2)$ theory following
Seiberg and Witten \cite{SW-I}.  $N=2$ has powerful Ward identities
which, together with some physical input, lead to interesting
conclusions like monopole condensation and - after $N = 2$ is softly
broken to $N = 1$ - to confinement due to condensation of the massless
monopoles. The chain of reasoning is long and elaborate.  For the
theories under consideration, when the $\theta$-angle is set to zero,
the BPS bound is given by
\begin{eqnarray}
M &\geq& \sqrt{2}\,\vert Z \vert\,,\nonumber \\
Z &=& a\,(n_e + {i\over \alpha} n_m)\,,\quad{\rm with}\quad 
\alpha = {g^2\over 4\pi}\,. \nonumber
\end{eqnarray}
The form of $Z$ is invariant under $n_e \leftrightarrow n_m$
accompanied by $\alpha\leftrightarrow 1/\alpha$ and $a \leftrightarrow
a/\alpha$. This may be regarded as some evidence for the
Montonen-Olive conjecture of electromagnetic duality, although we have
seen that $Z$ gets an interesting renormalization in $N=2$
theories. While the conjecture in its original form makes sense in
$N=4$ theories, for $N=2$, we can still have an $SL(2, Z)$
transformation acting on the parameter $\tau = \theta/2\pi +
i4\pi/g^2$ which embodies a weaker form of Montonen-Olive
conjecture. By combining these transformations with global symmetries
and the requirement of positivity of kinetic energy, Seiberg and
Witten were able to formulate a procedure for determining the exact
form of the low-energy theory as will be described below.

\subsection{Parametrization of the Moduli Space}

The classical action for the $N=2$ supersymmetric Yang-Mills theory
(\ref{N=2}) containts the scalar potential
$$
V={1\over 2g^2}{\rm Tr}\left([\phi^\dagger,\phi]^2\right)\,.
$$
Therefore, the Higgs vacuum is defined by $[\phi,\phi^\dagger]=0$,
which implies that $\phi$ takes values in the Cartan subalgebra of the
gauge group: $\phi= \phi_i H^i$. Thus, generically, the gauge group
$G$ is broken to the subgroup $H$ which is generated by elements from
the Cartan subalgebra. Elements in $G/H$ do not leave the Higgs vacuum
invariant, but being gauge transformations, they relate physically
equivalent vacua. On the other hand, once a given basis for the Cartan
subalgebra is chosen, then different vacuum values of $\phi$, within
this subalgebra, correspond to different physical theories. Thus these
degrees of freedom in $\phi$ ({\it i.e.}, the $\phi_i$'s) parametrize
the space of physically inequivalent vacua, or the moduli space of the
theory.  The dimension of this moduli space is equal to the rank $r$
of the gauge group $G$. However, this parametrization of the moduli
space is not the desired one as there is still some residual gauge
invariance: The coset $G/H$ contains elements which, while not leaving
the vacuum invariant, do not take $\phi$ out of the Cartan
subalgebra. These transformations are precisely the Weyl 
reflections. Therefore, the correct parametrization of the moduli
space is given not by $\phi$, but by Weyl invariant functions
constructed out of it.

The Weyl invariants are obtained from the characteristic equation, 
$$
\det (\lambda - \phi ) = 0\,.
$$
Since Weyl reflections act on $\phi$ by conjugation, $\det (\lambda
-\phi)$ is invariant. Hence, if we expand this quantity as a
polynomial in powers of $\lambda$, then the coefficients are Weyl
invariant quantities. In the following, we express these coefficients
as simple functions of $\phi$. Let $a_1,\cdots , a_N$ denote the roots
of the characteristic equation, or the eigenvalues of $\phi$. If the
gauge group is $SU(N)$ (or $SO(N)$), we have ${\rm Tr}\phi=\sum
a_i=0$. For generic values of $a_i$, the gauge group is broken to
$U(1)^r$. When some of the eigenvalues coincide, the unbroken group
jumps from $U(1)^r$ to something at least as big as $U(1)^{r-1}\times
SU(2)$. In terms of $a_i$, the characteristic polynomial takes the
form: 
$$
\lambda^N + \lambda^{N-2} \sum_{i<j} a_ia_j - \lambda^{N-3}
\sum_{i<j<k} a_ia_ja_k + \cdots + (-1)^N \prod_{i=1}^N a_i = 0\,.
$$
For $SU(2)$, $\phi={1\over 2}a\sigma_3$ and it can be easily checked
that the desired Weyl invariant is $u={\rm Tr}(\phi^2)={1\over2}
a^2$. For $SU(3)$, the coefficients in the characteristic polynomial
are  $a_1a_2+a_1a_3+a_2a_3$ and $a_1 a_2 a_3$, where $a_1, a_2, a_3$
are the eigenvalues of $\phi$. Using $({\rm Tr}\phi)^2=0$, and $({\rm 
Tr}\phi)^3=0$, one can easily write the Weyl invariants as    
$$
\begin{array}{ll}
u=\frac{1}{2}{\rm Tr}(\phi^2)=-(a_1a_2+a_1a_3+a_2a_3)\,,  \\
v=-\frac{1}{3}{\rm Tr}(\phi^3)=a_1 a_2 a_3\,.
\end{array}
$$
In general, for $SU(N)$ similar formulae can be worked out.  The
coefficients of the characteristic polynomial are the ``Chern" classes
of $\phi$, 
$$
\det(\lambda-\phi)=\lambda^N -\lambda^{N-1} c_{1}(\phi)+
\lambda^{N-2}c_{2}(\phi)+\cdots +(-1)^{j} \lambda^{N-j} c_{j} (\phi )+
\cdots + (-1)^N c_N(\phi)= 0\,.
$$
The coefficients $c_j(\phi)$ can be easily determined from the following
formal expansion,
\begin{eqnarray}
\det(\lambda -\phi) &=& \lambda^N \det (1-\phi/\lambda) =
\lambda^N e^{{\rm Tr\,ln}(1-\phi/\lambda)} 
=\lambda^N {\rm exp}\left(-\sum_{n=1}^\infty \frac{{\rm Tr}\,(\phi^n)}
{n\lambda^n}\right) \nonumber \\
&=& \lambda^N - \frac{1}{2}{\rm Tr}(\phi^2)\lambda^{N-2} -\frac{1}{3}
{\rm Tr}(\phi^3) \lambda^{N-3}+\cdots\,. \nonumber
\end{eqnarray} 
Note that the series expansion for ${\rm ln} (1-\phi/\lambda)$ makes
sense only for $\lambda >> \phi$. Therefore, in the above expansion
only terms with positive powers of $\lambda$ are relevant and they
provide all the Weyl invariant quantities we need for a gauge
invariant parametrization of the moduli space.

In the above, we have treated $\phi$ classically. In quantum field
theory, we parametrize the moduli space by the vacuum expectation
values of the corresponding classical Weyl invariants. For example,
the moduli space of the $SU(2)\,\,N=2$ supersymmetric Yang-Mills
theory is parametrized by $u=<{\rm Tr}(\phi^2)>$, which at the
classical limit, reduces to $a^2/2$.
 
\subsection{R-Symmetry and its Breaking}

The $N$-extended supersymmetry algebra (\ref{susy-noC}) is invariant
under global $U(N)$ rotations of the $N$ supercharges. Therefore, a
supersymmetric theory should also exhibit such a global symmetry,
usually referred to as R-symmetry. The action of this global symmetry
on the supercharges can be easily translated into a transformation of
the superspace variables. For example, for $N=1$, the $U(1)$
R-symmetry acts on the supercharge as $Q\rightarrow
e^{-i\alpha}Q$. From this we can obtain its action on the superspace
coordinates as $\theta\rightarrow e^{i\alpha}\theta$ and
$\bar\theta\rightarrow e^{-i\alpha}\bar\theta$ (see
(\ref{superspace-Q})).  For $N=2$, we have a $U(1)$ which acts as
above on the $N=2$ superspace coordinates $\theta^I, \bar\theta_I$,
along with an $SU(2)$ R-symmetry which rotates the index $I$ of the
supercharges. In order to keep the supersymmetric Lagrangian invariant
under these transformations, we have to assign appropriate
transformation properties to various superfields. From this, we can
easily obtain the behaviour of the component fields under
R-symmetry. Below, we will describe this in a little more detail.

\noindent\underline{Action on the $N=2$ Vector Multiplet:}
As discussed in the previous section, the $N=2$ vector multiplet
contains a vector field $A_\mu$, two Weyl spinors $\lambda, \psi$ and a
scalar $\phi$, all transforming in the adjoint representation of the
gauge group $G$. These components can be arranged as
$$
\begin{array}{ccc}
 & A_\mu & \\
 \lambda & & \psi \\
 & \phi &
\end{array}\,.
$$
The $SU(2)_R$ transformation acts on the rows in the above diagram and
rotates the fermions $(\lambda,\psi)$ into each other while keeping
$A_\mu$ and $\phi$ invariant. In the
$N=1$ formalism, this multiplet decomposes into a vector superfield
$V(A_\mu, \lambda)$, and a chiral superfield $\Phi(\phi,\psi)$.    
Therefore, the only part of $SU(2)_R$ which remains manifest in the
$N=1$ language is a $U(1)_J$ subgroup which does not mix $\lambda$ and
$\psi$. This subgroup is generated by $\sigma_3$ and acts as 
$(\lambda,\psi)\rightarrow(e^{i\alpha}\lambda,e^{-i\alpha}\psi)$. 
The action of the $U(1)_{\cal R}$ on the $N=1$ superfields was
discussed in the previous section. Below, we summarize these
transformations:
\begin{equation}
\begin{array}{lll}
U(1)_{\cal R} &: \Phi\rightarrow e^{2i\alpha}\Phi(e^{-i\alpha}\theta),
&V\rightarrow V(e^{-i\alpha}\theta)\,,
\\ \\
U(1)_J &: \Phi \rightarrow \Phi (e^{-i\alpha}\theta), &
V\rightarrow V(e^{-i\alpha}\theta)\,.
\end{array}
\label{Rb1}
\end{equation}

\noindent\underline{Action on the $N=2$ Hypermultiplet:}
In $N=2$, the matter fields appear in hypermultiplets, each containing
two complex scalars $(q,\widetilde q^\dagger)$ and two Weyl fermions
$(\psi_q,\widetilde\psi^\dagger_q)$. All these components transform in
the same (usually the fundamental) representation of the gauge group.
The components of a hypermultiplet can be arranged as 
$$
\begin{array}{ccc}
 & \psi_q & \\
 q & & \widetilde q^\dagger \\
 & \widetilde\psi_q^\dagger &
\end{array}\,. 
$$
Again, $SU(2)_R$ acts on the rows and thus rotates $q,\widetilde
q^\dagger$. In the $N=1$ language, the hypermultiplet is decomposed in
terms of chiral multiplets $Q(q,\psi_q)$, and $\widetilde Q(\widetilde
q,\widetilde\psi_q)$, which carry dual gauge quantum numbers. In this
decomposition, only the $U(1)_J$ subgroup of $SU(2)_R$ is
manifest. Moreover, in the $N=1$ decomposition, a hypermultiplet
interacting with a vector multiplet gives rise to a superpotential
term
$$
{\cal W} = \sqrt{2} \widetilde Q \Phi Q.
$$
Since ${\cal W}$ should carry two units of $U(1)_{\cal R}$ charge,
$Q$ and $\widetilde Q$ are neutral. We summarize these transformations
below: 
\begin{equation}
\begin{array}{lll}
U(1)_{\cal R}: &Q \rightarrow Q(e^{-i\alpha}\theta)\,,
& \widetilde Q \rightarrow \widetilde Q (e^{-i\alpha}\theta)\,, \\ \\
U(1)_J : &Q\rightarrow e^{i\alpha} Q(e^{-i\alpha}\theta)\,,
&\widetilde Q\rightarrow e^{i\alpha}\widetilde Q(e^{-i\alpha}\theta)\,. 
\end{array}
\label{Rb2}
\end{equation}
For later convenience, we list below the transformations of all
component fields under $U(1)_{\cal R}$ and $U(1)_J$:
$$
\begin{array}{lccc} 
U(1)_{\cal R}:\qquad & \phi &\rightarrow & e^{2i\alpha}\phi\,, \\
&(\psi,\lambda)&\rightarrow & e^{i\alpha}(\psi,\lambda)\,, \\
&(\psi_q,\widetilde\psi_q)&\rightarrow &e^{-i\alpha} 
(\psi_q,\widetilde\psi_q)\,, \\
&(A_\mu, q, \widetilde q)&\rightarrow & (A_\mu, q, \widetilde q)\,.
\\ \\
U(1)_J: &(\lambda, q)&\rightarrow & e^{i\alpha}(\lambda, q)\,,\\
&(\psi,\widetilde q^\dagger)&\rightarrow & e^{-i\alpha}
(\psi,\widetilde q^\dagger)\,,\\
&(\phi, A_\mu, \psi_q,\widetilde\psi_q^\dagger)&\rightarrow &
(\phi, A_\mu, \psi_q,\widetilde\psi_q^\dagger)\,.
\end{array}
$$
Note that we can combine the two-component spinors $\lambda$ and
$\bar\psi$ into a four-component Dirac spinor $\psi_D$ (see
(\ref{sp-comp})). The spinor $\psi_D$ transforms as $e^{i\alpha}
\psi_D$ under $U(1)_J$ and as $e^{i\alpha\gamma_5}\psi_D$ under
$U(1)_{\cal R}$. Similarly, $U(1)_{\cal R}$ acts as a chiral $U(1)$ on
the Dirac spinor constructed out of $\psi_q$ and $\widetilde
\psi_q^\dagger$\,, though now with the opposite charge. Thus
$U(1)_{\cal R}$ is a chiral symmetry and is, therefore, broken by  
a chiral anomaly as will be discussed below. 

\noindent\underline{Breaking of R-Symmetries:} 
Classically, our theory has the full global $SU(2)_R\times U(1)_{\cal
R}$ as a symmetry group. However, quantum mechanically, $U(1)_{\cal
R}$ is broken to a discrete subgroup due to anomalies. This subgroup
can be easily determined from elementary instanton considerations (the
more direct method will be described in the next subsection). To
compute the anomaly for the gauge group $SU(N_c)$, note that by the
index theorem, in the presence of an instanton, there is one zero-mode
for each left moving fermion in the fundamental or antifundamental
representaion and $2N_c$ zero-modes for each left-handed fermion in
the adjoint representation. A correlation function in this theory
involves integrations over the fermionic collective coordinates
corresponding to these zero-modes. For a correlator to be non-zero, it
should contain enough fermion insertions to soak the zero-modes. 
Hence, in the presence of $N_f$ flavours, the first non-vanishing
correlator is 
\begin{equation}
G=\langle\lambda(x_1)\cdots\lambda(x_{2N_c})\psi(y_1)\cdots\psi(y_{2N_c})
\psi_q(z_1)\cdots\psi_q(z_{N_f})\widetilde\psi_q(u_1)\cdots
\widetilde\psi_q(u_{N_f})\rangle \,.
\label{G}
\end{equation}
Under the $U(1)_{\cal R}$, $G$ transforms as
$$
G \rightarrow e^{i\alpha(4N_c-2N_f)}\,G\,.
$$
Hence $U(1)_{\cal R}$ is broken to the discrete group 
${\bf Z}_{4N_c-2N_f}$. In the following, we focus on the pure
Yang-Mills theory so that $N_f=0$. In this case, $U(1)_{\cal
  R}\rightarrow {\bf Z}_{4N_c}$ and is represented by $e^{2\pi
  i\alpha}$, where  $\alpha=n/4N_c\,, n=1,\cdots 4N_c$. Thus the
global symmetry group is $SU(2)_R\times {\bf Z}_{4N_c}$. However, note
that the centre of $SU(2)_R$, which acts as $(\lambda,\psi)\rightarrow
e^{i\pi}(\lambda,\psi)$, is also contained in ${\bf Z}_{4N_c}$
(corresponding to $n=2N_c$). Hence, the true symmetry group is  
$$
\left(SU(2)_R\times {\bf Z}_{4N_c}\right)/{\bf Z}_2\,.
$$

This surviving R-symmetry is broken further by the Higgs vacuum. 
The field $\phi^2$ has charge 4 under ${\bf Z}_{4N_c}$ and transforms
to $e^{2\pi i n/N_c}\phi^2$, which is invariant only for
$n=N_c,2N_c,3N_c,4N_c$. Therefore, if the vacuum is characterized by
non-zero $\phi^2$, then ${\bf Z}_{4N_c}$ is broken down to ${\bf
  Z}_4$. This is the situation for the $SU(2)$ gauge theory ($N_c=2$). 
In this case, all elements which do not keep $\phi^2$ invariant, act
as a ${\bf Z_2}: \phi^2\rightarrow -\phi^2$. Therefore, the final
R-symmetry group for the $SU(2)$ gauge theory is $(SU(2)_R\times 
{\bf Z_4})/{\bf Z_2}$. 

For the gauge group $SU(3)$, we have ${\bf Z_{4N_c}}={\bf Z_{12}}$.
The vacuum is parametrized by $\phi^2$ and $\phi^3$, which, under
${\bf Z}_{4N_c}$, transform to $e^{2\pi i n/N_c}\phi^2$ and $e^{3\pi i
n/N_c}\phi^3$ respectively. Invariance of $\phi^2$ requires
$n=N_c,2N_c,3N_c,4N_c$, breaking ${\bf Z_{12}}$ to ${\bf
Z_{4}}$. Invariance of $\phi^3$ picks $n=2N_c, 4N_c$. Thus, for the
$SU(3)$ gauge theory, ${\bf Z}_{12}$ is broken to ${\bf Z}_2$. The
broken ${\bf Z}_6$ subgroup acts non-tivially on the moduli space. For
$SU(4)$ and higher gauge groups, no subgroup of ${\bf Z}_{4N_c}$
survives.

\subsection{Low-Energy Effective Action for $N=2$ Gauge Theory}

Let us consider the $N=2$ supersymmetric Yang Mills theory based on a
group $G$ of rank $r$ which is spontaneously broken by a non-zero
$\langle\phi\rangle$. Far from the points where two eigenvalues of
$\langle\phi\rangle$ coincide, the only massless fields are the vector
supermultiplets associated with the unbroken subgroup $U(1)^{r}$ of
$G$. At sufficiently low energies, non of the massive states will
appear as physical states and an effective description of the theory
can be given in terms of the massless fields alone. In principle, this
low-energy theory can be obtained by integrating out all the massive
modes as well as massless modes above a low-energy cutoff. In
practice, this procedure is not easy to implement. Seiberg and Witten
formulated an indirect procedure for determining the exact low-energy
theory as we will in the remaining part of this section.

As discussed in the previous section, the effective action (the part
with at most two derivatives) is fully determined by a prepotential
${\cal F}$ which is only a function of $r$ massless vector
supermultiplets.  In the $N=1$ language, the corresponding Lagrangian
takes the form (with $A_i$ denoting chiral superfields):
$$
\frac{1}{4\pi} \, {\rm Im} \, \left [ \int \, d^{4}\theta \,
\frac{\partial {\cal F}}{\partial A^{i}} \bar{A}^{i} +
\int \, d^{2} \theta \, \frac{1}{2}
\frac{\partial^{2}{\cal F}}{\partial A_{i} \partial A_{j}}
W^{\alpha i} W^j_\alpha
\right ]\,.
$$
In terms of ${\cal F}(A)$, the  K\"{a}hler potential is given by
$K={\rm Im}\left(\bar{A}_{i}\partial{\cal F}(A)/\partial A_{i}\right)$. 
If $a_i$ denotes the scalar component of the chiral superfield $A_i$,
then the metric on the space of fields, and therefore, the one on the
space of Higgs vacuua, is given by
$$
ds^2 = g_{i\bar{j}}\,da^i d\bar{a}^j =
{\rm Im}\,\frac{d^2{\cal F}}{\partial a_i \partial a_j}
\, da^i d\bar{a}^j \,.
$$
Note that, by virtue of $N=2$ supersymmetry, this metric is the same
as the generalized gauge coupling which appears in front of the
$F_{\mu\nu}^iF^{\mu\nu i}$ term in the above Lagrangian. This
relationship is not modified in the quantum theory.   

In the following, we will be mainly interested in the $SU(2)$ gauge
group spontaneously broken to $U(1)$. In this case the effective
low-energy Lagrangian takes the form
\begin{equation}
{\cal L}_{eff}=
\frac{1}{4\pi} \, {\rm Im} \, \left [ \int \, d^{4}\theta \,
\frac{\partial {\cal F}(A)}{\partial A} \bar{A} +
\int \, d^{2} \theta \, \frac{1}{2}
\frac{\partial^{2}{\cal F}}{\partial A^{2}} W^{\alpha} W_{\alpha}
\right ]\,, 
\label{EA}
\end{equation}
and the metric on the field space is given by
\begin{equation}
ds^2= {\rm Im}(\tau)\, da d\bar a\,,\quad {\rm where},\quad
\tau (a) = \frac{\partial^{2} {\cal F}}{\partial a^{2}}.
\label{metric}
\end{equation}
Our aim in this subsection is to determine the perturbative form of
${\cal F}$ following \cite{Seiberg88}. The determination of the exact
form of the prepotential, including non-perturbative corrections, is
the subject of the later subsections. 

To determine the one-loop perturbative correction to ${\cal F}$, we
start from the exact microscopic theory (\ref{N=2}) with gauge group
$SU(2)$. This theory is asymptotically free and therefore at high
energies one can perform reliable perturbative calculations. As
indicated in the previous subsection, the $U(1)_{\cal R}$ symmetry of
this theory is broken by the standard chiral anomaly. Thus we have 
$$
\partial_\mu J^\mu_5= -\frac{N_c}{8\pi^2}F_{\mu\nu}\widetilde
F^{\mu\nu} \,.
$$
This implies that, to one-loop, under a $U(1)_{\cal R}$
transformation, the effective Lagrangian changes by 
\begin{equation}
\delta {\cal L}_{eff}= - \frac{\alpha N_c}{8\pi^2} F\widetilde F\,.
\label{anom}
\end{equation}
If the theory also contains $N_f$ fermions in the fundamental
representation, then in the above expression, $N_c$ has to be replaced
by $N_c-N_f/2$. Note that since $(32\pi^2)^{-1}\int F\widetilde F$ is
an integer, the anomaly breaks $U(1)_{\cal R}$ to ${\bf Z}_{4N_c}$ (or
to ${\bf Z}_{4N_c-2N_f}$ in the presence of matter). The same result
was obtained in the previous subsection from a consideration of
fermion zero modes in an instanton background.  Moreover, the anomaly
can be regarded as causing a shift in the $\theta$-angle. In the
following, we use this freedom to set $\theta$ to zero by an
appropriate chiral rotation of the fermions.
 
The one-loop form of ${\cal F}$ can be determined from the requirement
that under a $U(1)_{\cal R}$ transformation, the low-energy effective
action (\ref{EA}) change as in (\ref{anom}). The variation $\delta
{\cal L}_{eff}$ could come only from terms in ${\cal L}_{eff}$ which
are quadratic in $F_{\mu\nu}$. Writing only the relevent terms from
the Lagrangian (\ref{EA}), this means
$$ 
\frac{1}{16\pi}{\rm Im}\left[{\cal F}''(e^{2i\alpha}A)(-FF + iF
\widetilde F)\right] = \frac{1}{16\pi}{\rm Im}\left[{\cal F}''(A)
(-FF+iF\widetilde F)\right]-\frac{\alpha N_c}{8\pi^2}F\widetilde F\,.
$$ 
The form of the prepotential is therefore restricted by 
$$
{\cal F}''(e^{2i\alpha}A)={\cal F}''(A) - \frac{2\alpha N_c}{\pi}\,,
$$
or, for infinitesimal $\alpha$, by
$$
\frac{\partial^3 {\cal F}}{\partial A^3}=\frac{N_c}{\pi}\frac{i}{A}\,.
$$
The above equation can be easily integrated to give the one-loop
expression for the prepotential as 
\begin{equation}
{\cal F}_{1-loop}(A)=\frac{i}{2\pi} A^2 {\rm ln} \frac{A^2}{\Lambda^2}\,.
\label{oneloop}
\end{equation}
Here, $\Lambda$ is a fixed dynamically generated scale like
$\Lambda_{QCD}$.  It is known that, due to $N=2$ supersymmetry, the
above one-loop expression for the prepotential does not receive higher
order perturbative corrections. This is related to the fact that in
this theory the perturbative $\beta$-function is only a one-loop
effect. Thus, in this theory, the one-loop nature of the perturbative
$\beta$-function is consistent with the well known one-loop nature of
the anomaly. This is not the case in $N=1$ theories where the anomaly
is, of course, still a one-loop effect but not the $\beta$-function
(see, for example, \cite{SV1, SV2}).

Although, the prepotential, as given in (\ref{oneloop}), is exact in
perturbation theory, it does receive non-perturbative corrections due
to instanton effects as argued in \cite{Seiberg88}. The general form
of these corrections can be fixed as follows: First, it is clear that
a correction to ${\cal F}$ coming from a configuration of instanton
number $k$ should be proportional to the $k$-instanton factor
$exp(-8\pi^2k/g^2)$ (since the prepotential is a holomorphic function,
it cannot receive corrections from anti-instanton configurations). 
Using the one-loop $\beta$-function of the theory given by
$\beta(g)=-g^3/4\pi^2$, the $k$-instanton factor can be written as 
$$
e^{-8\pi^2k/g^2} = \left(\frac{\Lambda}{a}\right)^{4k}\,.
$$
Furthermore, following the approach of Seiberg in \cite{Seiberg-n,
Seiberg-poh}, we note that the broken $U(1)_{\cal R}$ symmetry is
restored if we assign a charge of $2$ to $\Lambda$. With this
modification, the prepotential should transform under $U(1)_{\cal R}$
as a field of charge $4$, without a non-homogeneous term. This implies
that the $k$-instanton correction should also be proportional to
$A^2$. Putting these together, the prepotential including generic
non-perturbative corrections can be written as
$$
{\cal F}=\frac{i}{2\pi} A^2 {\rm ln} \frac{A^2}{\Lambda^2}
+\sum_{k=1}^\infty\,{\cal F}_k\left(\frac{\Lambda}{A}\right)^{4k}A^2\,.
$$
The coefficients ${\cal F}_k$ are not field dependent due to the fact
that in a supersymmetric theory, instantons contribute to the path
integral only through zero-modes\cite{DD}. The coefficient ${\cal
F}_1$ was calculated in \cite{Seiberg88} and found to be non-zero. The
determination of the exact form of ${\cal F}$ is the subject of the
work of Seiberg and Witten. 

The one-loop expression for ${\cal F}$ can also be obtained from
integrating the expression for the $\beta$-function, choosing $A$ as
one of the integration limits: In the classical low-energy theory,
which is obtained by simply dropping the massive fields from the
microscopic Lagrangian (and not integrating over them), the
prepotential is given by ${\cal F}(A)=\frac{1}{2} \tau_{cl}A^2$,
where, $\tau_{cl}=\frac{\theta}{2\pi}+i\frac{4\pi}{g^2}$. In the
quantum theory, for large $a$ (which is the {\it vev} of the scalar
somponent of $A$), asymptotic freedom takes over and the theory is
weakly coupled. Therefore, in this limit, a good approximation to the
quantum behaviour of the theory can be obtained by replacing $g$ by
the corresponding running couplings $g(a)$ (we have set $\theta$ to
zero by a chiral redefinition of the fermion fields).  In the limit of
large $a$, $g(a)$ and hence ${\cal F}(a)$ can be obtained in
perturbation theory by integrating the $\beta$-function.

For other gauge groups we note that for every root $\alpha$, ($\alpha
> 0$), we have massive fields $W^{\alpha}, W^{- \alpha}$. For
simply-laced groups, $\alpha^{2} = 2$, and
$$
{\cal F} = \frac{i}{4\pi} \, \sum_{\alpha > 0} \, (\vec{\alpha} \cdot
\vec{a})^{2} \, {\rm ln}\, \frac{(\vec{\alpha}\cdot
\vec{a})^{2}}{\Lambda^{2}}\,.
$$
The coefficient of the logarithm follows from the one-loop
$\beta$-function and it also leads to the unbroken global
$U(1)_{R}$-symmetry.

The one-loop expression for ${\cal F}$, coupled with the fact that a
well defined theory should have a positive kinetic energy term, leads
to very interesting consequences.  For large $|a|$, using
(\ref{oneloop}), we can calculate $\tau(a)=\frac{i}{\pi}({\rm
ln}\,\frac{a^2}{\Lambda^2}+ 3)$. This is a multi-valued function,
though the metric on the field space given by ${\rm Im}\tau$ is single
valued. However, since ${\rm Im}\tau(a)$ is a harmonic function, it
cannot have a global minimum. Thus, if it is globally defined, it
cannot be positive everywhere. Therefore, the positivity of the
kinetic energy requires that $\tau(a)$ is defined only locally. This
means that in the region of the complex plane where $\tau(a)$ becomes
negative, one needs a different description of the theory. In the next
subsection we describe how these equivalent descriptions could be
obtained.

\subsection{Duality}
To find the duality transformations which relate different
descriptions of the same theory, we consider the gauge field terms in
the bosonic part of the action. Working in the Minkowski space with
conventions $(F_{\mu\nu})^2=-(\widetilde F_{\mu \nu})^2$ and
$\widetilde{\widetilde F}=-F$, these
terms can be written as 
$$
\frac{1}{32 \pi} \, {\rm Im}\, \int \, \tau (a) (F+i\widetilde F)^2 =
\frac{1}{16 \pi} \, {\rm Im} \, \int \, \tau (a) (F^{2}+i\widetilde
FF) \,.
$$
Now we regard $F$ as an independent field and implement the Bianchi 
identity $dF = 0$ by introducing a Lagrange multiplier vector field
$V_{D}$. To fix the Lagrange multiplier term, $U(1)\subset SU(2)$ is
normalized such that all $SO(3)$ fields have integer charges. Then, as
discussed in section 1, all matter fields in the fundamental
representation of $SU(2)$ will have half-integer charges. With this 
convention, a magnetic monopole satisfies $\epsilon^{0\mu\nu\rho}
\partial_{\mu}F_{\nu\rho}=8\pi\delta^{(3)}(x)$. The Lagrange
multiplier term can now be constructed by coupling $V_D$ to a
monopole:
$$
\frac{1}{8\pi}\,\int\,V_{D \mu}\epsilon^{\mu\nu\rho\sigma}
\partial_{\nu}F_{\rho\sigma}=\frac{1}{8\pi}\,\int\,\widetilde F_{D}F = 
\frac{1}{16\pi}\,{\rm Re}\,\int\,(\widetilde
F_{D}-iF_{D})(F+i\widetilde F)\,,
$$
where, $F_{D\mu\nu}=\partial_\mu V_{D\nu}-\partial_\nu V_{D\mu}$. 
Adding this to the gauge field action and integrating over $F$,
gives the dual theory
$$
\frac{1}{32\pi}\,{\rm Im}\int \left(-\frac{1}{\tau}\right)
(F_D +i\widetilde F_D)^2=\frac{1}{16\pi}\,{\rm Im}\int\left(-\frac{1}{\tau}
\right) (F_D^2 +i\widetilde F_D F_D)\,.
$$
The dualization can also be performed in an $N=1$ supersymmetric
language. In this case, the relevant term in the action is 
$$
\frac{1}{8 \pi} \, {\rm Im} \, \int \, d^{2}\theta
\tau (A) W^{2}\,.
$$
The Bianchi identity is replaced by ${\rm Im}{\cal D}W = 0$. This can
be implemented by introducing a vector superfield $V_{D}$ and the
corresponding Lagrange multiplier term becomes
$$
\frac{1}{4\pi}\,{\rm Im}\,\int d^4x d^4\theta\, V_{D}{\cal D}W 
=\frac{1}{4 \pi}\,{\rm Re}\,\int d^4x d^4\theta\, i{\cal D}V_D W
=-\frac{1}{4\pi}\,{\rm Im}\,\int d^4x d^2\theta\, W_{D} W\,.
$$
Adding this to the action and integrating out $W$, gives the dual
action 
$$
\frac{1}{8\pi}\,{\rm Im}\,\int d^2\theta \left(-\frac{1}{\tau(A)}
W^2_D \right)\,. 
$$
Thus, the effect of the duality transformation is to replace a gauge
field which couples to electric charges by a dual gauge field which
couples to magnetic charges, and at the same time, transform the
gauge coupling as  
\begin{equation}
\tau\rightarrow\tau_D = -\,\frac{1}{\tau}\,.
\label{tautauD}
\end{equation}
We recognize this as the electric-magnetic duality of section 1.
Also, note that the action is invariant under the replacement
$\tau\rightarrow\tau+1$. This transformation, along with the one in
(\ref{tautauD}), generates the $SL(2,{\rm Z})$ group which, therefore,
is the full duality group of our theory. This group acts on $\tau$ as 
\begin{equation}
\tau\rightarrow{\frac{a\tau + b}{c\tau + d}}\,,
\label{tauSLtwoZ}
\end{equation}
where, $ad-bc=1$ and $a,b,c,d \in {\rm Z}$. 

$N=2$ supersymmetry relates the dual description of the gauge fields
to a dual description for the scalar fields. To see this, let us
introduce $h(A)=\partial{\cal F}/\partial A$. In terms of this,
$\tau(A)=\partial h(A)/\partial A$ and the scalar kinetic energy term
becomes ${\rm Im}\int d^4\theta\,h(A)\bar A$. For the dual theory
corresponding to $(\ref{tautauD})$, let us introduce the corresponding
dual variables $A_D,{\cal F}_D,h_D(A_D)$ and $\tau_D$. Then, equation
(\ref{tautauD}) implies that $A_D=h=\partial{\cal F}/\partial A$ and 
$h_D = -A$. Under this transformation, the scalar kinetic energy term
retains its form,
$$
{\rm Im}\int d^4\theta\,h(A)\bar A=
{\rm Im}\int d^4\theta\,h_D(A_D)\bar A_D\,. 
$$
In the following, we use the notations $A_D$ and $h(A)$
interchangeably. 

We now consider the effect of the full $SL(2,Z)$ group on $A$ and
${\cal F}$, or equivalently on $A$ and $A_D=\partial{\cal F}/\partial
A$. The transformation (\ref{tauSLtwoZ}) implies that
\begin{equation}
\left(\begin{array}{c}A_D\\ A\end{array}\right)\rightarrow
\left(\begin{array}{cc} a & b \\ c & d \end{array}\right)
\,\left(\begin{array}{c}A_D\\ A \end{array}\right) \,.
\label{ASLtwoZ}
\end{equation} 
Note that, in general, we could also add a constant column matrix $C$
to the right-hand side of the above equation. However, for non-zero
$C$, the BPS mass formula for the theory in the absence of matter is
not invariant under the above transformation. Thus in this case we
should set $C=0$. However, in the presence of matter fields, the BPS
mass formula is modified and a non-zero $C$ is allowed. This case will
be discussed in the last section. On the space of the scalar fields,
the transformations $\tau\rightarrow -1/\tau$ and $\tau\rightarrow
\tau+1$ are implemented by the matrices
\begin{equation}  
S=\left(\begin{array}{cc}0&1\\-1&0\end{array}\right),\quad
T=\left(\begin{array}{cc}1&1\\0&1\end{array}\right).
\label{ST}
\end{equation}
These matrices generate the $SL(2,{\rm Z})$ group.

The transformation of ${\cal F}$ can be easily obtained from
(\ref{ASLtwoZ}), or equivanlently from 
\begin{eqnarray}
A'_D &=& a A_D + b A\,, \nonumber\\
A'   &=& c A_D + d A\,. \nonumber
\end{eqnarray}
The first equation can be integrated with respect to $A'$ by using the
second equation and the result is 
$$
{\cal F}'=\frac{1}{2}\beta\delta A^2 +\frac{1}{2}\alpha\gamma A^2_D
+\beta\gamma A A_D + {\cal F}\,.
$$

Let us now come back to the metric on the moduli space as given by 
(\ref{metric}). In terms of the variable $a_D$, this takes the form
$$
ds^2={\rm Im}\,da_D d\bar{a}=-\frac{i}{2}(da_D d\bar{a}-da d\bar{a}_D)\,,
$$
which is invariant under the $SL(2,{\rm Z})$ transformations. To make
the description more precise, note that the moduli space ${\cal M}$ is
a complex one dimensional manifold and let $u$ be a holomorphic
coordinate on this manifold. Finally, we will identify $u$ as 
$\langle{\rm Tr}\phi^2\rangle$. Let $a$ and $a_D$ be the coordinates
on a space $X \cong {\rm C}^2$ on which we can choose a symplectic form
$\omega={\rm Im}\,da_D\wedge d\bar{a}$. The functions
$(a_{D}(u),a(u))$ give a map $f$ from ${\cal M}$ to $X$. In other
words, they determine a section of $X$ regarded as an $SL(2,{\rm Z})$
bundle over ${\cal M}$. In terms of the coordinate $u$, the metric on
the moduli space takes the form 
$$
ds^2 ={\rm Im}\,\frac{da_D}{du}\,\frac{d \bar{a}}{d\bar{u}}\,
dud\bar{u}=-\frac{i}{2}\left(\frac{da_{D}}{du}\,\frac{d\bar{a}}{d\bar{u}}
-\frac{d\bar{a}_D}{d\bar{u}}\,\frac{da}{du}\right)dud\bar{u}\,.
$$
This is the pull-back of the K\"{a}hler metric associated with the
symplectic form $\omega$ and is, therefore, manifestly $SL(2,{\rm Z})$
invariant. Choosing $u=a$, we get back to the original formula. Note
that for arbitrary $(a_D(u),a(u))$, the metric is not
positive. However, the solution we will obtain later determines these
functions such that the metric is always positive.  

\noindent\underline{Higher Dimensions:}
{}For a group of rank $r$, the metric takes the form
$$
ds^2 ={\rm Im}\,\frac{\partial^2{\cal F}}{\partial a_i\partial a_j}\, 
da_i d\bar{a}_j\,.
$$
Introduce a space $X \simeq {\bf C}^{2r}$ with coordinates $a^{i},
a_{Dj}=\partial {\cal F}/\partial a^{j}$. Endow $X$ with the symplectic
form $\omega=\frac{i}{2}\,\sum_{i}(da^i\wedge d\bar{a}_{Di}-da_{Di} 
\wedge d\bar{a}^i)$ and the holomorphic 2-form $\omega_h =\sum_{i} \, 
da^{i} \wedge da_{Di}$. Choose $u^s\, (s=1,...,r)$ as coordinates on
the moduli space ${\cal M}$. Then, $a^i(u),a_{Dj}(u)$ give a map $f$
from ${\cal M}$ to $X$ such that $f^{*}(\omega_{h}) = 0$ (this
condition ensures that $a_{Dj}=\partial {\cal F}/\partial a_{j}$). 
Thus, the metric takes the form 
$$
ds^{2} = {\rm Im} \, \frac{\partial a_{Di}}{\partial u^{r}} \,
\frac{\partial \bar{a}^{i}}{\partial \bar{u}^{s}} \,
du^{r} d\bar{u}^{s}\,.
$$
This is the metric associated with $f^*(\omega)$ and is therefore
invariant under $Sp(2r,{\rm R})$: If we write $v^T =(a_D, a )$, then
the metric is invariant under $v\rightarrow Mv + c$, where, $M \in
Sp(2r,{\bf R})$. In the absence of matter, $c=0$ and moreover, only
the $Sp(2r,{\bf Z})$ part survives.

If the moduli space ${\cal M}$ has a non-trivial structure, then, on
being taken around a close loop, the vector $v^T=(a_D, a)$ will get
transformed by an element of the monodromy group. As  will be shown
later, the monodromy group is a subgroup of $SL(2,{\rm Z})$, and its
determination is essential for solving the problem. 

\subsection{Dyon Masses}

As described in sections 1 and 2, for the microspcopic $SU(2)$ theory,
the BPS bound is given by $M \geq \sqrt{2} |Z|$, where,
$$
Z=a(n_e + \tau_{cl}\, n_m)\,, 
$$
and $\tau_{cl}=\frac{\theta}{2\pi}+i\,\frac{4\pi}{g^2}$. All states
that saturate this bound fall in a short multiplet of the $N=2$
algebra. In a more general $N=2$ theory, like the low-energy effective
action we have been studying, this formula is slightly modified. 
Suppose, the theory contains matter fields in hypermultiplets. When
$a\ne 0$, these fields, which in the $N=1$ formalism are described by
chiral fields $M, \widetilde M$, become massive. If a hypermultiplet
carries charge $n_e$, then its coupling to the chiral field $A$ is
uniquely fixed by $N=2$ supersymmetry as 
$$
\sqrt{2}\,n_e \, A M \widetilde M\,.
$$
From this we can easily see that for such a state, $Z=a\,n_e$. On the
other hand, if we consider a magnetic monopole carrying $n_m$ units of
magnetic charge, then a manipulation very similar to the one in
subsection 1.6 leads to the BPS bound $\sqrt{2}|n_m\,a_D|$ or,
equivalently, $Z=n_m\,a_D$. In general, one can compute this formula
for dyons and obtain 
\begin{equation}
Z = a n_e + a_D n_m\,.
\label{ZaaD}
\end{equation}
If the gauge group is of rank $r$, then we have
$$
Z = a^i n_{e,i} + h_i(a) n^i_m =a\cdot n_e +a_D \cdot n_m\,.
$$
Since masses are physically observable, the mass formula should be
invariant under the monodromies on the moduli space. Therefore, if
$v=(a_D, a)^T$ transforms to $M v$ (with $M\in Sp(2r,{\rm Z}) )$, then
the vector $w= (n_m, n_e)$ should transform to $w M^{-1}$. Note that
this once again requires $M$ to be an integral matrix. 

Under the action of the monodromy, the coupling matrix $\tau_{ij}=
\partial a_Di/\partial a_j$ transforms to $(A\tau+B)(c\tau+D)^{-1}$,
where, $M=\left(\begin{array}{cc}A&B\\C&D\end{array}\right)$. This is
very similar to the transformation of the period matrix of a genus $r$
Riemann surface under the action of the monodromy group on the moduli
space of genus $r$ surfaces. This makes it reasonably appealing to
identify our $r$-dimensional moduli space of vacua ${\cal M}$ with the
moduli space of genus $r$ Riemann surfaces. The variables $a_i,a_{Dj}$
are then related to the periods of these surfaces and can be
calculated. Moreover, ${\rm Im}\,\tau$ is always positive definite,
which resolves the problem of geting a negative kinetic term. To check
this hypothesis for our $SU(2)$ theory, we first have to determine the
monodromy structure on the moduli space ${\cal M}$. This is the
subject of the next subsection.
